% =============================================================
% Chapter 1: Overview — Who's Who in the Age of AI
% Book: AI Figures
% =============================================================

\chapter{AI 群星闪耀时：一部关于人的技术史}
\label{ch:overview}

% Chapter overview: Why a biographical encyclopedia of AI matters,
% methodology, scope, and reading guide.

技术史的本质是人的历史。

每一次技术革命的背后，都站着一群做出关键决策的人——
他们选择了某个研究方向而非另一个，创办了某家公司而非留在学术界，
签下了某笔投资而非另一笔，推动了某项政策而非放任自流。
这些个体决策的总和，塑造了我们今天看到的AI格局。

2012年，Alex Krizhevsky、Ilya Sutskever和Geoffrey Hinton
用一个叫AlexNet的卷积神经网络在ImageNet上取得了突破性成绩，
深度学习从边缘理论变成了主流范式。
2017年，Google的八位研究员发表了``Attention Is All You Need''，
Transformer架构从此成为几乎所有大语言模型的基础——
而这八个人后来分别创办或加入了六家不同的AI公司，
他们的职业选择本身就是一部微型AI产业史。
2022年底，Sam Altman决定将ChatGPT以免费网页的形式发布，
这个看似简单的产品决策引爆了人类历史上最大规模的技术创业浪潮。

本书是LLM Observatory系列的第三部专著，
与姊妹篇\textit{Emergence}（技术篇）和
\textit{The AI Startup Landscape}（产业篇）共同构成
对AI时代的三维透视。如果说\textit{Emergence}回答的是
``这些模型是怎么训练出来的？''，
\textit{The AI Startup Landscape}回答的是
``这些公司是怎么运作和竞争的？''，
那么本书回答的核心问题是：
\textbf{这些事是谁做的？他们为什么做出了那样的选择？
他们的背景、经历和人际网络如何塑造了整个AI生态？}


% =============================================================
\section{为什么需要一部AI人物百科？}
\label{sec:overview-why}
% =============================================================

\subsection{人际网络即技术路线图}

AI领域的一个显著特征是：
\textbf{关键人物的流动方向，往往比任何技术路线图都更能预测
产业的下一步走向}。

当Dario Amodei和Daniela Amodei带领十余名研究员
从OpenAI出走创办Anthropic时（2021年），
这不仅仅是一次人事变动——它预示了AI安全将成为一个独立的
商业赛道，Anthropic后来成长为估值超过3800亿美元的巨头。
当Noam Shazeer离开Google创办Character.AI，
又被Google以27亿美元的许可协议``收回''时（2024年），
这揭示了大公司通过``反向收购''获取人才的新模式。
当梁文锋决定将幻方量化的GPU集群投入基础模型研发时（2023年），
一个量化基金悄然变成了震动全球的DeepSeek。

理解这些人物——他们的教育背景、职业轨迹、师承关系、
合作网络和个人信念——是理解整个AI时代的必要条件。

\subsection{本书的独特价值}

市面上不缺关于AI技术的教科书，也不缺关于AI产业的商业分析。
但系统性地、百科全书式地梳理AI时代\textbf{所有}重要人物
的详细履历、贡献和人际关联，这在中英文出版物中都是首次尝试。

本书覆盖了\textbf{超过700位}在AI时代（2012--2026）
发挥关键作用的人物，包括：

\begin{itemize}
  \item \textbf{深度学习先驱}——三位Turing Award得主、
    Transformer八位作者、LSTM发明者等奠基人物
  \item \textbf{AI实验室领袖}——OpenAI、Anthropic、Google DeepMind、
    Meta AI、xAI、Mistral等核心团队
  \item \textbf{中国AI生态}——DeepSeek、智谱、月之暗面、MiniMax、
    百川、阶跃星辰等``六小虎''，以及BAT和华为的AI团队
  \item \textbf{芯片与基础设施}——NVIDIA Jensen Huang、
    AMD Lisa Su、TSMC Morris Chang等算力巨头
  \item \textbf{投资人与资本}——a16z、Sequoia、红杉中国、
    SoftBank等塑造AI资本格局的关键人物
  \item \textbf{创业者}——从Cursor到Perplexity，从Midjourney到Suno，
    覆盖编程、搜索、创意、企业服务等全赛道
  \item \textbf{学术界}——Stanford、Berkeley、CMU、MIT、
    清华、北大等顶尖实验室的核心研究者
  \item \textbf{政策、伦理与新星}——监管者、批评者、
    开源社区英雄和35岁以下的明日之星
\end{itemize}


% =============================================================
\section{方法论与数据来源}
\label{sec:overview-method}
% =============================================================

本书的人物信息来源包括：

\begin{enumerate}
  \item \textbf{公开履历}：LinkedIn、个人网站、大学主页、
    公司团队页面
  \item \textbf{学术记录}：Google Scholar、Semantic Scholar、
    DBLP的发表和引用数据
  \item \textbf{新闻报道}：TechCrunch、The Information、
    Bloomberg、36氪、量子位等科技媒体
  \item \textbf{公司信息}：Crunchbase、PitchBook融资数据，
    SEC/港交所公开披露文件
  \item \textbf{社交媒体}：Twitter/X、微博等平台上的
    公开发言和互动
  \item \textbf{播客与访谈}：Lex Fridman Podcast、
    Acquired、20VC等长形式访谈
\end{enumerate}

需要特别说明的是：本书聚焦于\textbf{公开可获取的信息}。
我们不揣测私人动机，不转述未经证实的传闻，
对于有争议的事件（如OpenAI 2023年11月的董事会危机），
我们尽量呈现多方视角并标注信息来源的可靠程度。

\begin{warnbox}[关于时效性的说明]
AI领域的人事变动极为频繁——平均每周都有重要的人才流动消息。
本书的信息截止日期为\textbf{2026年3月}。
对于在截止日期后发生的变动，读者可参考本书的在线更新版本。
某些估值、融资金额和职位信息可能在出版后发生变化。
\end{warnbox}


% =============================================================
\section{本书结构与阅读指南}
\label{sec:overview-structure}
% =============================================================

本书共14章（含附录），按以下逻辑组织：

\begin{corebox}[章节结构]
\begin{enumerate}
  \item \textbf{第1章}（本章）：全景概览与方法论
  \item \textbf{第2--7章}：按组织/实验室分类的西方AI人物
    \begin{itemize}
      \item 第2章：深度学习先驱与教父级人物
      \item 第3章：OpenAI与GPT革命
      \item 第4章：Anthropic与AI安全运动
      \item 第5章：Google DeepMind
      \item 第6章：Meta AI与开源运动
      \item 第7章：xAI、Mistral、Cohere及其他西方实验室
    \end{itemize}
  \item \textbf{第8章}：中国AI生态（最长章节，100+人物）
  \item \textbf{第9章}：芯片与基础设施巨头
  \item \textbf{第10章}：AI投资人与资本
  \item \textbf{第11--12章}：AI创业者（消费/企业/垂直/机器人）
  \item \textbf{第13章}：学术研究者
  \item \textbf{第14章}：政策、伦理、媒体与新星
  \item \textbf{附录}：人物总索引与统计
\end{enumerate}
\end{corebox}

\subsection{推荐阅读路径}

\subsubsection{技术背景读者}
从第2章（先驱）开始，按顺序阅读至第7章，
然后跳到第13章（学术界）了解师承关系，
最后回到第8章（中国）和第14章（新星）。

\subsubsection{投资/商业背景读者}
先读第10章（投资人），
然后读第3章（OpenAI）和第4章（Anthropic）了解头部公司，
接着读第11--12章（创业者）和第9章（基础设施），
最后读第8章（中国）了解东方格局。

\subsubsection{政策/社会视角读者}
从第14章（政策与伦理）开始，
然后读第2章（先驱，关注AI安全思想源流），
第4章（Anthropic与安全运动），
最后读第8章（中国）了解地缘政治维度。


% =============================================================
\section{与姊妹篇的关系}
\label{sec:overview-siblings}
% =============================================================

本书与系列中其他两部专著的关系可以用三个维度来理解：

\begin{keybox}[LLM Observatory 三部曲]
\begin{itemize}
  \item \textbf{Emergence}（涌现）——\textit{技术维度}：
    Transformer、预训练、后训练、推理优化、Agent……
    回答``模型是怎么做出来的？''
  \item \textbf{The AI Startup Landscape}（AI创业全景）——
    \textit{产业维度}：基础设施层、模型层、平台层、应用层……
    回答``产业是怎么运作的？''
  \item \textbf{AI Figures}（本书）——\textit{人物维度}：
    先驱、研究者、创业者、投资人、政策制定者……
    回答``这些事是谁做的？''
\end{itemize}
三本书的交叉阅读将提供对AI时代最完整的理解。
本书在相关处会引用姊妹篇的具体章节，读者可按需跳转。
\end{keybox}


% =============================================================
\section{一个时代的群像}
\label{sec:overview-closing}
% =============================================================

Stefan Zweig在《人类群星闪耀时》中写道：
``一个真正具有历史意义的时刻——一个人类的群星闪耀时刻——
出现以前，必然会有漫长的岁月无谓地流逝。''

AI时代的群星闪耀时，恰恰发生在我们这一代人的眼前。
从2012年的AlexNet到2026年的多模态Agent，
不过短短十四年，却可能是人类文明最密集、
最深远的技术变革之一。

参与这场变革的不是抽象的``技术力量''或``市场趋势''，
而是一个个具体的人——有的出身名校、
有的自学成才，有的追求利润、有的信仰开源，
有的乐观到近乎天真、有的悲观到呼吁暂停。
他们的故事，就是AI时代的故事。

翻开下一页，让我们从三位Turing Award得主开始，
走进这个群星闪耀的时代。
